\section{Worked kernels: batch norm}
\label{sec:bn}

Batch norm is four forward kernels in \code{kernels/teeny-kernels/src/nn/norm/batchnorm.rs},
and which one runs depends on the op and the mode:

{\small
\begin{xltabular}{\linewidth}{>{\raggedright\arraybackslash}p{2.7cm}>{\raggedright\arraybackslash}p{4.6cm}>{\raggedright\arraybackslash}X}
\toprule
\textbf{Op, mode} & \textbf{Kernel(s)} & \textbf{Spec in the lowered DAG} \\
\midrule
\endhead
BN2d, inference & \code{batch_norm_2d_nchw_forward_inference} & Yes, rank 4 via \code{span} (PR \#103). \\
BN2d, training & the same inference kernel & Yes, but the kernel normalises with \emph{running} statistics (\issue{teenygrad-fq0l}). \\
BN1d/3d, inference & \code{batch_norm_forward_inference} (\code{#[kernel]}) & No. The kernel declares nothing, and \code{make_float_kernel!} sets \code{tile_spec: None}. \\
BN1d/3d, training & \code{batch_norm_stats_forward} $\to$ \code{batch_norm_normalize_forward} & No. Both kernels derive specs, but the lowering attaches \code{None} (\issue{teenygrad-1nr.18.7}). \\
\bottomrule
\end{xltabular}
}

\subsection{Inference, NCHW: a broadcast map over a spanned axis}

\begin{lstlisting}
#[tiled_kernel]
pub fn batch_norm_2d_nchw_forward_inference<T: Triton, D: Float, const BLOCK_HW: i32>(
    #[tile(extent = B, dim = Y)]
    #[tile(extent = C)]
    #[tile(block = BLOCK_HW, extent = HW, span = 2)]
    x: In<Tile<T, D>>,
    #[tile(extent = B, dim = Y)]
    #[tile(extent = C)]
    #[tile(block = BLOCK_HW, extent = HW, span = 2)]
    y: Out<Tile<T, D>>,
    #[tile(extent = C)] weight: In<Tile<T, D>>,          // broadcast operands: a subset of
    #[tile(extent = C)] bias: In<Tile<T, D>>,            // the output's axes
    #[tile(extent = C)] running_mean: In<Tile<T, D>>,
    #[tile(extent = C)] running_var: In<Tile<T, D>>,
    B: i32, C: i32, HW: i32, eps: f32,
) {
    let rstd = T::rsqrt(running_var.tensor + T::cast(T::full(&[BLOCK_HW], eps), ..));
    let y_tile = weight.tensor * (x.tensor - running_mean.tensor) * rstd + bias.tensor;
    T::store(y.tensor, y_tile, x.mask, &[], None, None);
}
\end{lstlisting}

\begin{itemize}
\item \textbf{The prelude.} \code{B} is decoded from \code{program_id(Y)}, and
  \code{C} and \code{HW} from \code{program_id(X)}. \code{HW} is the one blocked
  range, and \code{x}'s offsets are \code{__tile_range + tile_c * HW + tile_b * (C * HW)}.
\item \textbf{The four per-channel operands} declare only \code{C}, which makes them
  broadcast operands. Each loads one element at \code{arange(0, 1) + tile_c} and
  \code{broadcast_to}s \code{[BLOCK_HW]}. The body is plain arithmetic on equal-shaped
  tiles.
\item \textbf{The spec.} \code{x}/\code{y} are rank~4 with
  \code{untiled_dims: ["B", "C"]} and one binding \code{{dims: [2, 3], block:
  "BLOCK_HW", extent: "HW"}}. The four operands are rank~1 with no binding and
  \code{untiled_dims: ["C"]}.
\item \textbf{The grid.} \code{[B (Y), C, HW (BLOCK_HW)]}.
\end{itemize}

\textbf{Resolution.} An output tile \code{[1, 1, 1, 128]} gives \code{resolved =
{HW: 1*128}}, so \code{x = [None, None, 1, 128]}. Every per-channel operand is
\code{[None]}: \emph{all} \code{C} channels, though a program reads exactly one. This
is the conv2d problem again. \code{C} is grid-decoded, one per program, but carries
no block, so nothing seeds it.

\subsection{Training, step 1: \texttt{batch\_norm\_stats\_forward} (reduce-only)}

\begin{lstlisting}
#[tiled_kernel]
#[tile_loop(trip_count = [N, BLOCK_N])]
#[tile_carry(acc_sum = [BLOCK_N], acc_sum_sq = [BLOCK_N])]    // see the gotcha below
#[tile_reduce_pass(over = N, read = x, into = sum,
    acc = T::sum(x, None, true), finish = sum * (1/N), store = mean)]
#[tile_reduce_pass(over = N, read = x, into = sum_sq,
    acc = T::sum(x * x, None, true), finish = sum_sq * (1/N) - sum * sum, merge)]
pub fn batch_norm_stats_forward<T: Triton, D: Float, const BLOCK_N: i32>(
    #[tile(block = BLOCK_N, extent = N, reduce, walk)]     // [N, C]: N outer, walked, reduced
    #[tile(block = 1, extent = C)]
    x: In<Tile<T, D>>,
    #[tile(block = 1, extent = C)] mean: Out<Tile<T, D>>,
    #[tile(block = 1, extent = C)] rstd: Out<Tile<T, D>>,
    running_mean_ptr: InOut<T::Pointer<D>>,                // untagged: not in the spec
    running_var_ptr: InOut<T::Pointer<D>>,
    N: i32, C: i32, eps: f32, momentum: f32,
) { /* rstd = rsqrt(var + eps); EMA update of the running stats */ }
\end{lstlisting}

\textbf{The spec.} \code{x} is rank~2. \code{N} is \code{walk}, so it has no binding
and goes to \code{untiled_dims: ["N"]}. \code{C} gets a binding with
\code{block_const "1"}, and \code{reduction_axis} is \code{Some(0)}. \code{mean} and
\code{rstd} are rank~1, each with a \code{C} binding of block \code{"1"}. The grid is
\code{[C]}.

\textbf{Resolution.} A \code{mean} tile \code{[k]} gives \code{resolved = {C: k}},
so \code{x = [None, k]}: $k$ channels of \emph{every} row. That is exactly right. A
statistic needs its whole reduction domain, and ``\code{None} = full extent'' is here
the truth rather than a fallback. Two operands are invisible to any cost model,
because \code{running_mean_ptr}/\code{running_var_ptr} are untagged.

\begin{note}[Gotcha: the declared carries are not the real ones]
\code{#[tile_carry(acc_sum = [BLOCK_N], acc_sum_sq = [BLOCK_N])]} predates the
conversion to \code{#[tile_reduce_pass]}. The generated body accumulates into
\emph{scalar} carries named \code{sum} and \code{sum_sq} (\code{zeros(&[1])}), and no
variable called \code{acc_sum} exists. The \code{loop_spec} this kernel reports is
therefore stale metadata. Nothing consumes \code{loop_spec} yet, which is why
nothing failed.
\end{note}

\subsection{Training, step 2: \texttt{batch\_norm\_normalize\_forward} (map pass)}

\begin{lstlisting}
#[tiled_kernel]
#[tile_grid(order = [C])]                                   // one program per channel; N is walked
pub fn batch_norm_normalize_forward<T: Triton, D: Float, const BLOCK_N: i32>(
    #[tile(block = BLOCK_N, extent = N, walk)] #[tile(block = 1, extent = C)] x: In<Tile<T, D>>,
    #[tile(block = BLOCK_N, extent = N, walk)] #[tile(block = 1, extent = C)] y: Out<Tile<T, D>>,
    #[tile(block = 1, extent = C)] weight: In<Tile<T, D>>,
    #[tile(block = 1, extent = C)] bias: In<Tile<T, D>>,
    #[tile(block = 1, extent = C)] mean: In<Tile<T, D>>,
    #[tile(block = 1, extent = C)] rstd: In<Tile<T, D>>,
    N: i32, C: i32,
) {
    let mean_b = T::broadcast_to(mean.tensor, &[BLOCK_N]);   // ... rstd_b, gamma, beta likewise
    gamma * (x - mean_b) * rstd_b + beta                       // one iteration of the map
}
\end{lstlisting}

There is no reduce pass. A \code{walk} axis on its own selects the multi-pass path,
whose map loop is the whole kernel. The prelude loads the four statistics once, as
\code{[1]} tiles at \code{tile_c}. The map loop walks \code{N} in \code{BLOCK_N}
steps with stride \code{C}, loads \code{x}, evaluates the body, and stores.

\textbf{Resolution.} A \code{y} tile \code{[?, k]} gives \code{{C: k}}, so
\code{x = [None, k]} and every statistic is \code{[k]}. Compare the NCHW inference
kernel: there the same ``one channel per program'' relationship is declared
\emph{without} a block, and the operands come back at full \code{C}. Two spellings
of one fact propagate differently. This is the strongest argument for the rule
proposed in \S\ref{sec:conv-overestimate}: a gridded axis is a block of 1, whether
or not the author wrote \code{block = 1}.

\subsection{Why NCHW training BN does not fit today's vocabulary}
\label{sec:bn-train}

The stats kernel's input is a 2-D \code{[N, C]} view. That is a real BatchNorm1d
tensor, and also what BatchNorm2d/3d look like in a channels-last layout, where
\code{B, H, W} are adjacent and flatten into one \code{N}. In NCHW the reduction for
a channel runs over \code{B} \emph{and} \code{H, W}, two runs of dims separated by
\code{C}. None of the three ways of saying that is available:

\begin{itemize}
\item \textbf{\code{reduce} on two axes:} \code{reduction_axis} is one index, and the
  macro rejects a second \code{reduce}.
\item \textbf{\code{span} over the reduced dims:} \code{span} cannot combine with
  \code{reduce}, and \code{B} and \code{HW} are not adjacent anyway.
\item \textbf{\code{walk} over both:} the pass generator walks the first \code{walk}
  axis it finds.
\end{itemize}

The cheapest unlock is to let \code{reduction_axis} become a list. With that, an NCHW
stats kernel could declare \code{B} and \code{HW} both \code{reduce} (\code{HW} with
\code{span = 2}), and \code{resolve_inputs} would need no change, because reduced dims
are unbound and resolve to full extent regardless. A layout-transposed
(channels-last) path is the alternative. Either one is a prerequisite for
\issue{teenygrad-fq0l}.
